\newpage
\section{Indépendance}

Deux évènements sont indépendants si le comportement de l'un n'affecte pas le comportement de l'autre.

\ex{
Alice tire une carte au hasard dans un jeu de 52 cartes.\\
Bob lance un dé à 6 faces au même moment.\\
Quelle est la probabilité qu'Alice tire une reine et que Bob obtienne un 5?\\

\vspace{2cm}
\noindent
Charles lance 2 fois de suite un dé. Il n'est pas affecté émotionnellement par les résultats de ses dés.\\
Quelle est la probabilité qu'il obtienne 2 six ?
}
\vspace{2cm}
\begin{tcolorbox}[colback=cyan!5, colframe=cyan!70!blue, title=Définition de l'indépendance de deux évènements]
\defn{Deux évènements $A$ et $B$ sont \textbf{indépendants} ssi
\begin{equation}
    \mathbb{P}(A\cap B) = \mathbb{P}(A) \mathbb{P}(B)
\end{equation}}
\end{tcolorbox}

\ex{
\noindent
Charles lance à nouveau 2 fois de suite un dé. Il n'est toujours pas affecté émotionnellement par les résultats de ses dés.\\
Quelle est la probabilité que la somme de ses résultats soit égale à 8 ?\\

\vspace{1cm}
\noindent
Quelle est la probabilité que le résultat du premier dé soit pair ?\\

\vspace{1cm}
\noindent



Quelle est la probabilité que la somme de ses résultats soit égale à 8  et que le résultat du premier dé soit pair ?\\


}